% Paper-ready validity, ethics, and reproducibility text for AdverSim.
% This section is additive and does not change reported results.

\section{Validity, Ethics, and Reproducibility}

\subsection{Threats to Validity}

\textbf{Construct validity.}
The structural realism score $R(S)$ measures ATT\&CK consistency, scenario
structure, and transition coherence. It should not be interpreted as a complete
proxy for real-world incident realism. The human validation experiment shows
that human overall-realism ratings do not strongly positively align with
$R(S)$, so the metric is best framed as a structural consistency measure that
requires human realism checks.

The SOC simulation is parameterized and should be interpreted as simulation
behavior rather than an empirical production-SOC study. Sensitivity analysis is
therefore used to distinguish invariant findings from model artifacts. Sigma
replay provides defensive grounding by applying generated rule skeletons to
labelled events, but its coverage depends on event availability, label quality,
logsource mapping, and the intentionally transparent keyword matcher.

\textbf{Internal validity.}
The validate-and-requery loop reduces invalid JSON, schema failures, and
inactive ATT\&CK IDs, but it can also bias accepted outputs toward scenarios
that satisfy schema and matrix constraints. Human ratings were collected from
five independent technical reviewers: two software test engineers, one IT
specialist, one master's student in cybersecurity, and one master's student in
computer science. This supports independent human realism checking, but it
should not be described as a panel composed entirely of SOC, DFIR, or
threat-intelligence experts.

\textbf{External validity.}
The corpus covers selected attack types, environments, difficulties, and
ATT\&CK Enterprise techniques. Results may not transfer directly to operational
incidents, non-Enterprise ATT\&CK matrices, highly specialized industrial
control environments, or organizations with different logging maturity. The real
baseline comparison uses public CALDERA Stockpile and Atomic Red Team metadata;
these are useful non-author-constructed baselines but are not complete models of
all real adversary behavior.

\textbf{Statistical conclusion validity.}
Confidence intervals, null models, held-out refits, agreement statistics, and
multiple-testing adjustments are used where the available data support them.
Remaining risks include limited human-reviewer count, aggregate-only legacy
mutation outputs, and correlation estimates that depend on the selected sample
and metric definition.

\subsection{Ethics and Safety}

AdverSim is intended for defensive cyber-range research, teaching,
reproducible evaluation, and detection engineering. The repository generates
structured scenario descriptions and Sigma-compatible detection skeletons. It
does not execute malware, exploit payloads, Atomic Red Team tests, CALDERA
agents, or live adversary emulation. Public baseline scripts parse metadata
only.

Raw defensive logs, raw human rating sheets, reviewer identities, API keys, and
local secrets are excluded from version control. Environment files and key-like
artifacts are ignored, raw annotation directories are ignored, and tracked
outputs contain only aggregate or anonymized results. Out-of-scope uses include
unauthorized intrusion, malware deployment, credential theft, exploit execution,
or evasion testing against systems without authorization.

\subsection{Artifact Availability and Reproducibility}

The repository separates API-dependent generation from offline
reproducibility. The offline path runs without API keys or raw logs: it checks
tracked result artifacts, regenerates paper figures from aggregate outputs,
validates figure hashes through a manifest, runs the smoke pipeline, and scans
for large files or secret-like values. Fresh LLM generation, multi-model
replication reruns, and some public-baseline refreshes are documented separately
because they require API access or network availability.

\begin{table}[t]
\centering
\caption{Methodological components and their role in strengthening validity.}
\begin{tabular}{p{0.23\linewidth}p{0.31\linewidth}p{0.31\linewidth}}
\hline
Component & Contribution & Limitation \\
\hline
Validate-and-requery & Rejects malformed JSON and invalid ATT\&CK IDs before corpus acceptance & Does not guarantee full tactical realism \\
Single ATT\&CK source of truth & Makes technique counts traceable to one matrix snapshot & Depends on selected ATT\&CK version \\
Cross-version robustness & Tests behavior under v13/v14/v15 matrix drift & Limited to Enterprise ATT\&CK \\
Statistical rigor & Adds confidence intervals, null models, and effect-size framing & Some legacy outputs remain aggregate-only \\
SOC sensitivity & Separates invariant findings from model artifacts & Not a production-SOC study \\
Sigma replay & Grounds detection claims in labelled events & Depends on labels and simple matching \\
Real baselines & Compares against public non-author-constructed metadata & Public corpora are heterogeneous \\
Human validation & De-biases $R(S)$ with blinded human ratings & Reviewers are technical, not all SOC/DFIR experts \\
Reproducibility CI & Regenerates figures and verifies artifacts without secrets & Fresh generation still requires API access \\
\hline
\end{tabular}
\end{table}

\subsection{Data and Model Card Summary}

\begin{table}[t]
\centering
\caption{AdverSim data and model card summary.}
\begin{tabular}{p{0.32\linewidth}p{0.56\linewidth}}
\hline
Field & Summary \\
\hline
Artifact & Synthetic cyber-range scenario corpus in JSONL format \\
Scenario structure & Narrative, environment, difficulty, attack type, stages, ATT\&CK techniques, and attack graph \\
Main ATT\&CK matrix & Enterprise v14.1 active-technique index \\
Robustness checks & Additional v13/v14/v15 validation audit \\
Human validation & 150 blinded scenarios rated by five independent reviewers \\
Reviewer composition & Two software test engineers, one IT specialist, one MSc cybersecurity student, one MSc computer science student \\
Raw sensitive files & API keys, raw logs, raw rating sheets, and reviewer identities excluded from git \\
Intended uses & Cyber-range scenario generation, defensive evaluation, teaching, and reproducible research \\
Out-of-scope uses & Unauthorized intrusion, exploit execution, malware deployment, credential theft, or unsanctioned evasion testing \\
\hline
\end{tabular}
\end{table}

\subsection{Limitations Statement}

AdverSim scenarios are synthetic and should not be interpreted as substitutes
for real incident reports or telemetry. The validation pipeline reduces
malformed outputs and invalid ATT\&CK IDs, but tactical realism remains partly a
human judgment. Accordingly, $R(S)$ is reported as a structural consistency
metric, while blinded human ratings provide an independent realism check. SOC
detection results are presented as simulation and sensitivity-analysis outputs,
not as measurements of analyst performance in a production SOC.
